% GENERATED by scripts/make_tables.py -- do not edit by hand.
\begin{table}[t]
    \centering
    \small
    \begin{tabular}{lrrrrr}
        \toprule
        framework & total & none & any & proper interval & bare dispersion only \\
        \midrule
\multicolumn{6}{l}{\textit{Panel A: eligible records}} \\
        Overall & 96 & 81 & 15 & 9 & 6 \\
        Anthropic & 33 & 28 & 5 & 1 & 4 \\
        OpenAI & 32 & 30 & 2 & 0 & 2 \\
        DeepMind & 21 & 17 & 4 & 4 & 0 \\
        Third-party & 10 & 6 & 4 & 4 & 0 \\
        \addlinespace
\multicolumn{6}{l}{\textit{Panel B: source documents contributing eligible records}} \\
        Overall & 21 & 13 & 8 & 5 & 3 \\
        Anthropic & 7 & 4 & 3 & 1 & 2 \\
        OpenAI & 5 & 4 & 1 & 0 & 1 \\
        DeepMind & 4 & 2 & 2 & 2 & 0 \\
        Third-party & 5 & 3 & 2 & 2 & 0 \\
        \bottomrule
    \end{tabular}
    \caption{Uncertainty reporting at two denominators. At record level, ``none''
    means that row reports no uncertainty; at source-document level, it means
    that none of the document's eligible records reports uncertainty. A document
    is counted under ``proper interval'' if at least one eligible record reports
    one; ``bare dispersion only'' means it reports uncertainty somewhere but no
    proper interval. These are descriptive counts for the fixed purposive frame,
    not estimates from a probability sample.}
    \label{tab:reporting-levels}
\end{table}
